% ─── Discussion ───────────────────────────────────────────────────────────────
\label{sec:discussion}

% The discussion follows the research-question structure: each question is
% addressed in full before the next, and the cross-cutting policy implications
% (carbon pricing, infrastructure risk) are gathered after the empirical
% questions they bear on. RQ1 (6.1, 6.2), RQ2 (6.3, 6.4), the coal-grid trajectory
% qualification that bears on both (6.5), policy implications of RQ1 (6.6 ETS2,
% 6.7 stranded assets), then RQ3 scope (6.8) and the limitations that bound them (6.9).

%── 6.1 Overview ───────────────────────────────────────────────────────────────
\subsection{Interpreting the Results in Light of the Research Questions}

This paper set out to answer three questions. Two are empirical and are taken up
here, namely under what conditions decarbonised
hydrogen is cost-competitive with heat pumps for residential heating across the
EU-27, the United Kingdom and Switzerland (RQ1), and how Switzerland can draw on
the wider European hydrogen build-out, including imports, to serve its own
net-zero buildings strategy (RQ2). The evidence yields negative conclusions for
both questions, though for reasons that differ from those the framing of the
questions anticipated. Hydrogen heating does not drive a building-heat transition
on its own economics (RQ1), and it offers limited strategic value for
Switzerland's net-zero buildings strategy despite that country's access to
European supply (RQ2). Hydrogen reaches between under 1 and about 8 per cent of useful heat at
the modelled end-date, and only under scenarios with substantial carbon pricing
(Figure~\ref{fig:fig6}). That ceiling is set by the peaking niche hydrogen can
win, by the investment economics of the plant that would serve it, and by a
base-load comparison that now points the same way.

Once equipment capital is
charged against machines sized to peak heat demand, once each carrier pays for
its own last-mile network, and once hydrogen pays for the seasonal store that
its winter-weighted load shape creates, the best available heat pump undercuts
the hydrogen boiler in 22 of 29 countries by 2050, at a median gap of
$+$\euro{}15.4/MWh (Figure~\ref{fig:fig5}). Five of the seven exceptions are
salt-cavern markets. Space heating draws about three-quarters of its annual
energy in the October-to-March half of the year while electrolytic production is
far flatter, so a hydrogen heating load has to be stored across the seasons, and
hydrogen's base-load case concentrates where the geology makes that store cheap. This is
the same endowment that decides the three peaking arenas, the building winter peak, the
district-heating stack and the power-sector peaker.

The arenas therefore close the same door for one physical reason, and their agreement
is worth less as corroboration than it may appear. All three price the hydrogen side
through the same storage parameter in the same functional form, so their country
rankings are correlated by construction, and finding that they line up is closer to a
property of the model than to three independent confirmations. What they do establish
is scope, since an operating-cost win in one arena does not carry to another and none
clears the capital-recovery hurdle. No hydrogen peaker recovers its capital from market
rent in any country or any scenario, on the assumed \euro{}60/MWh scarcity premium and
with no capacity payment, and what base-load lead survives in the seven rests on a
fiscal asymmetry between an untaxed molecule and taxed retail electricity, not on a
technology advantage.

The finding qualifies the framing embedded in REPowerEU \citep{REpowerEU2022} and in
several national hydrogen strategies, which implicitly position hydrogen as a near-term
substitute for natural gas in buildings. The analysis suggests instead that the
transition will be heat-pump-led, and that any material role for hydrogen in buildings
before 2050 would have to rest on targeted policy support justified by co-benefits.
Where the delivered-cost comparison still favours hydrogen it does so on a
tax-inclusive electricity price set against an untaxed molecule, so a strategy built on
that comparison rests on a policy choice a future government can reverse, in the only
countries where the comparison survives at all.

%── 6.2 Comparison with Prior Literature ───────────────────────────────────────
\subsection{Comparison with Prior Literature on Heating}

The results align with, but qualify, a maturing body of comparative work.
\citet{Berger2020} found that heat pumps dominate gas boilers on lifecycle cost
across most of Europe by 2030 under moderate carbon pricing. The present analysis
supports that finding on the 2050 horizon, where the best available heat pump
undercuts the gas boiler in 24 of 29 countries on the central carbon path, and
qualifies it in the near term, where the two stand at parity at the study median
in 2025 and the heat pump is ahead in 16 of 29. That 2050 count is the one result
here that moves with the carbon price, running from 19 of 29 on the low path to
all 29 on the high path, precisely because the gas boiler pays ETS2 and the heat
pump does not. The qualification traces to equipment capital, which is
charged here against a machine sized to peak heat demand and so does not shrink
with the seasonal coefficient of performance. Figure~\ref{fig:fig5}(b) ranks the 2050 levelised cost
of heat by country and technology under the central carbon and hydrogen
trajectories, plotting the best available heat pump, the hydrogen boiler and the
gas boiler for each.

\citet{Huckebrink2022} concluded that heat pumps outperform hydrogen boilers for German
residential heating, attributing the result to the coefficient-of-performance advantage.
The present model still does not reproduce that ordering for Germany, and the
discrepancy is worth setting out. Reading the German entry of that ranking, the best
available heat pump reaches \euro{}139.3/MWh by 2050 and the hydrogen boiler comes in
about \euro{}25/MWh below it, which places Germany among the seven markets in which
hydrogen still leads on base load. Three things put it there, and unusually cheap German
hydrogen is not among them. The first is that equipment capital is charged against a
machine sized to peak heat demand, which withdraws an implicit credit that a high
seasonal coefficient of performance would otherwise confer on the heat pump's capital
charge.

The second, and the larger, is the German treatment of retail electricity, in which
levies, taxes and network charges lift the household tariff far above the wholesale cost
of the kilowatt-hour, while hydrogen is costed delivered and untaxed and carries no
equivalent burden. The third is geological. Germany has salt caverns, so it holds its
seasonal hydrogen imbalance at about \euro{}1.5/kg, worth \euro{}12.1/MWh of useful
heat, against \euro{}28.4/MWh in a cavern-free market such as Spain, and it is that
endowment that keeps Germany on hydrogen's side of a line that most of the study area
now sits the other side of. The German case is therefore a statement about relative
fiscal treatment and storage geology. It makes no claim that the
coefficient-of-performance mechanism has ceased to operate. That mechanism is intact, and a levy-neutral comparison
would restore much of the ordering \citet{Huckebrink2022} reports.

The cost-stack decomposition of Figure~\ref{fig:fig7}, set out in full in
Appendix~\ref{sec:results:lcohanatomy}, bears on that literature in three ways. It
corroborates the emphasis on the running-cost penalty of combustion heating, since fuel
is the largest single component of all three stacks. It locates hydrogen's disadvantage
in the logistics of the molecule and not in the cost of producing it, since delivery and
seasonal storage take more than a quarter of the hydrogen stack. And it does not support
treating the heat pump as a capital-dominated technology set against fuel-dominated
rivals, because fuel is the largest item in the heat-pump stack too, at 54.2~per cent,
even though the heat pump carries the heaviest capital share of the three.

The scenario
comparison of \citet{Billerbeck2023Race}, whose model-based 'race' between
hydrogen and heat pumps for space and water heating finds heat pumps winning
across most configurations, and the techno-economic caveats catalogued by
\citet{Badakhsh2023Caveats} for green hydrogen in buildings, reach the same
verdict on the aggregate outcome. The results here agree on that outcome and
locate it more precisely, since the ceiling on hydrogen arises from the cost of
holding a winter-weighted load across the seasons, which acts on the base-load
comparison, the peaking dispatch and capital recovery alike, and which is lower in
the seven markets whose salt caverns can hold that store cheaply.

The emissions ranking is consistent with the life-cycle literature.
\citet{Slorach2021Net} showed that hydrogen boilers carry higher lifecycle emissions
than heat pumps even on low-carbon hydrogen, owing to upstream production and
distribution losses; the planetary-boundaries assessment of \citet{Weidner2023Planetary}
reaches a comparable verdict across a wider set of impact categories. The emissions
results here corroborate the ordering. Net Zero achieves the lowest near-term CO$_2$
trajectory, at 339~MtCO$_2$ by 2030, against 513~MtCO$_2$ for H2 Push, 532 for Stated
Policies and 548 for Current Policies, because deep renovation combined with grid
decarbonisation drives heat-pump emissions down faster than hydrogen production chains
can themselves decarbonise.

At the spatial margin, \citet{Maestre2022} found that hydrogen may be competitive in
specific high-density urban building segments in Spain yet is unlikely to achieve
mass-market adoption on economic grounds; the country-level results here place Spain
firmly on the heat-pump side of the base-load comparison, because it lacks the salt
caverns that halve the seasonal-storage charge, so any Spanish hydrogen case rests on
the dense urban segments Maestre and colleagues identify and not on the national cost
ordering. The contribution relative to this literature is threefold: a residential
heat-demand reconstruction built bottom-up at NUTS3 resolution, a merit-order test that
prices hydrogen on short-run marginal cost in a dispatch, not on levelised cost alone,
and an investment test that asks whether the peaker can recover its capital, which the
levelised comparisons do not pose.

%── 6.3 Switzerland (RQ2) ──────────────────────────────────────────────────────
\subsection{Switzerland's Low Hydrogen Demand for Heating, a Consequence of Absent Cavern Storage}

The Swiss case answers RQ2 directly, and every arena now answers it the same way. On
base-load levelised cost Switzerland sits on the heat pump's side of the line, at a 2050
heat-pump band of \euro{}122.4 to 131.1/MWh against \euro{}128.0/MWh for the hydrogen
boiler, so the best Swiss heat pump is about \euro{}5.6/MWh the cheaper. The reason is
the same one that orders the study area as a whole. Switzerland has no salt caverns, so
a Swiss hydrogen heating load is stored at the \euro{}3.0/kg rate, not the \euro{}1.5/kg
cavern rate, and that charge is what holds expensive, heavily taxed Swiss retail
electricity ahead of an imported molecule. The case against direct hydrogen heating
there is reinforced by the other two arenas and by the country's position in the supply
chain.

Hydrogen wins no Swiss peak in any of the four scenarios once the Swiss reservoir fleet
is credited as firm winter capacity, the Swiss hydrogen peaker costing \euro{}135 to
236/MWh-heat against a heat-pump peak of \euro{}62 to 90/MWh-heat. Without that credit the
uncoupled 29-country comparison prices the Swiss peak at \euro{}155/MWh-heat and hydrogen
does win it under H2~Push, which is why Switzerland is one of the 16 in that row; the
credit removes the win, and the endogenous arm is the better-specified test for a
country whose reservoir headroom the uncoupled series does not see. No Swiss peaker
recovers its capital on either arm; and Switzerland is cavern-poor and import-dependent, so
every molecule burned in a Swiss home is one bought on a competitive European market.
Switzerland illustrates a more general principle alongside this. Countries that combine
a clean electricity grid with high building-fabric standards weaken the \emph{emissions}
case for hydrogen in buildings relative to the EU-27+UK+CH average, so the question
that matters is why direct hydrogen heating is a poor strategic bet there in the first
place. France, whose
grid intensity of 80~gCO$_2$/kWh in 2025 falls to 8~gCO$_2$/kWh by 2050, and Sweden,
where it falls from 45 to 5~gCO$_2$/kWh \citep{EMBER2024electricity, EEA2023fit55},
share this characteristic. For these countries the policy priority is to accelerate
heat-pump installation and the supporting grid reinforcement.

\subsection{What Switzerland Can Draw from the European Build-Out}
\label{sec:discussion:swissstrategy} The finding that direct hydrogen heating is a poor
strategic bet for Switzerland does not exhaust RQ2, since it leaves open whether
Switzerland can draw on the European hydrogen build-out for ends other than heat. Two qualifications close the loop in that
direction. First, the low domestic demand for hydrogen heat is a strength rather than a
vulnerability for an import-dependent, cavern-poor system. A heat-pump-led buildings
sector spares Switzerland from competing for scarce imported molecules to warm its
homes, and frees its access to European supply for uses where the molecule carries a
higher value.

Second, those higher-value uses lie outside buildings. Swiss access to the emerging
European hydrogen network may create value as a consumer for hard-to-abate industry,
parts of heavy transport, and a firm-capacity role in the power system against the
country's winter-electricity adequacy gap, the very sectors in which hydrogen
substitutes for costly or absent alternatives and not for a mature heat pump. The
strategic answer to RQ2 is therefore that its leverage over the European build-out is
best realised away from the building stock, where the same access that does little for
residential heating may yet underwrite decarbonisation in the harder sectors; the
buildings-to-power link is taken up under RQ3 below, and the wider industrial and
transport uses are left to future work.

\subsection{Coal-Grid Markets and the Trajectory Argument}
\label{sec:discussion:coalgrid} The harder case is the set of Eastern European countries
whose grids are still carbon-intensive today, and here the argument turns on the
trajectory of grid carbon intensity rather than on its present level. In Poland, at 600~gCO$_2$/kWh in 2025
\citep{EMBER2024electricity}, heat pumps are already cost-competitive on levelised cost
yet carry a higher near-term operating carbon footprint than in the clean-grid
countries. That near-term penalty is transitional. The EEA projections have these grids
decarbonising substantially by 2040, at which point heat-pump carbon performance
converges with that of Western Europe \citep{EEA2023fit55}. The policy implication
follows directly from this convergence. Eastern European countries should prioritise the
gas-to-heat-pump switch now, while accelerating grid decarbonisation in parallel,
instead of treating present grid carbon intensity as a reason to defer electrification,
since the carbon penalty is temporary and the infrastructure choices made now are durable.

Figure~\ref{fig:fig12} provides the supporting evidence for this convergence, tracing
the grid-intensity trajectories of the 29 modelled countries from 2025 to 2050 with the
most carbon-intensive, the largest and the cleanest grids marked. Because heat-pump
operating emissions scale directly with grid intensity, the narrowing of the
inter-country band from a 2025 spread of roughly 600 to 45~gCO$_2$/kWh to about a
quarter of that width by 2040, and to an order of magnitude smaller by 2050, is what
removes the near-term carbon gap on which the
deferral argument rests, and quantifies the penalty that an Eastern-European heat pump
carries over a Western-European one in any given year as the vertical distance between
the trajectories.

\begin{figure}[H]
    \centering
    \caption{Electricity grid carbon-intensity trajectories, 2025 to 2050, on the exogenous
    paths of Section~\ref{sec:method:emissions}. The figure has no scenario dimension. It
    shows the median country, the tenth-to-ninetieth percentile band across the 29, and
    Poland, Germany and Sweden as the most carbon-intensive, the largest and the cleanest
    grids. These are inputs to the model and not results of it.}
    \includegraphics[width=0.95\textwidth]{paper/P19_grid_trajectories.pdf}
    \label{fig:fig12}
    \\[2pt]{\footnotesize\textit{Source: Authors' contribution.}}
\end{figure}

%── 6.6 Policy implication I: the structural role of ETS2 ──────────────────────
\subsection{Policy Implication I: The Structural Role of ETS2 Carbon Pricing}

The RQ1 finding carries two policy implications that the remainder of the
discussion draws out; the electricity merit-order and peaking results that inform
them are the substance of RQ3 and are developed in Section~\ref{sec:power}, so the
present subsections read them as implications of the heating result. The first implication concerns carbon pricing. The
second emissions trading system, ETS2, which extends carbon pricing to
buildings and road transport from 2027, does more to move the pathway than any other policy lever in the
pathway, and its incidence on the cost stack is quantified in
Appendix~\ref{sec:results:lcohanatomy} (Figure~\ref{fig:fig7}). The central claim of this
subsection is that carbon pricing accelerates the transition instead of determining its direction. Heat pumps already stand at parity with gas boilers on
levelised cost before ETS2 takes effect, at the study median in 2025 and ahead
outright in 16 of 29 countries, so the carbon price acts as an accelerant on an
ordering that is close to balanced without it.

The differential ETS2 burden is what tips the remaining markets. Each \euro{}50/tCO$_2$ of carbon price adds about \euro{}11/MWh of useful heat to the gas boiler, so the central path carries a 2030 median of about \euro{}27/MWh against a median \euro{}0.7/MWh on the heat pump. It closes the residual economic case for retaining gas boilers where
they are still ahead, and raises the opportunity cost of delayed electrification
everywhere else. The pass-through of the carbon price to the consumer
is a material source of uncertainty. The model samples pass-through between 75 and
100 per cent, a range consistent with the empirical estimates of
\citet{Fabra2009etspassthrough}; incomplete pass-through, for instance where
governments cap energy bills, would blunt the price signal and slow the
transition.

This carbon-pass-through mechanism, set out above, recurs at the
building peak, the district-heat stack and the power peak alike, and the country counts
across the three arenas bear it out. Table~\ref{tab:arenas} reports, for each of
the four scenarios (Current Policies, Stated Policies, Net Zero and H2 Push), the
count of the 29 countries in which hydrogen undercuts its fossil rival on operating
cost in three arenas, the building winter peak, the district-heating peak and the
power-sector peak. No row falls as carbon-pricing ambition rises, every row starts at
zero under Current Policies, and the two building-scale rows reach their widest count
under H2~Push, so no scenario without a carbon-price wedge on the fossil competitor
opens the niche at all. The even-handed power row does not rise past seven from Net
Zero on, because the carbon wedge that opens the niche is already at its full width
there.

%── 6.7 Policy implication II: stranded-asset and infrastructure risk ──────────
\subsection{Policy Implication II: Infrastructure and Stranded-Asset Risk in Buildings}

The second implication concerns infrastructure. If the Stated Policies or Net Zero
pathways materialise, the gas distribution
grids that serve residential buildings face steeply declining utilisation from
the 2030s, which raises the risk of a tariff spiral. Falling volumes force fixed
network costs onto a shrinking customer base through higher unit charges, the
higher charges accelerate disconnections, and the disconnections raise charges
further. This dynamic has been modelled for the United Kingdom by
\citet{Competition2021} and bears directly on Germany's municipal heat-planning
(\emph{W\"armeplanung}) mandate under the 2024 Building Energy Act.

Figure~\ref{fig:infrabill} (Appendix~\ref{sec:results:infrabill}) quantifies the scale
of the exposure. Under H2~Push the pathway builds a hydrogen distribution network and
still needs more electricity reinforcement than that network costs, because hydrogen
serves only the peak sliver while heat pumps carry the bulk of the load, so hydrogen
does not, in aggregate, substitute for the network investment that electrification
demands. The H2
Push pathway only partly forecloses the tariff spiral, since repurposing the gas grid to carry hydrogen preserves utilisation for the 9.3 per cent
peak niche that the merit order supports under that scenario
(Section~\ref{sec:results:meritorder}), and the missing-money result implies
that even that niche would require a capacity payment to be financed privately.

That 9.3 per cent is a merit-order and financing ceiling that the base-load comparison does not dispute, since delivered cost favours hydrogen only in a similarly narrow set, and the dispatch awards it a role in the seven cavern-endowed markets and in no
others. The infrastructure-bill comparison therefore sharpens the stranded-asset tension
instead of resolving it. The trade-off between stranding gas assets and building out
hydrogen infrastructure prematurely is not squarely resolved in current EU or national
strategy documents, and it warrants explicit policy analysis, the more so because the
pipeline-blending route that is sometimes invoked as a bridge is itself bounded by the
technical and regulatory blend limits reviewed by \citet{Erdener2022Review} and by the
distribution-grid material constraints assessed by \citet{Giehl2023}.

%── 6.8 Cross-Sector Synergies (RQ3 scope) ─────────────────────────────────────
\subsection{Cross-Sector Synergies and the Scope of RQ3} \label{sec:discussion:rq3}

Beyond the carbon-pricing and infrastructure implications above, a third question
follows. RQ3 asks what synergies hold between hydrogen in buildings and its use in
other sectors. This paper answers the buildings-to-power coupling in full, on the same
operating-cost basis as the rest of the analysis (Section~\ref{sec:power}), and
defers the wider synergies with industry and transport to subsequent work, for the
reasons set out in Section~\ref{sec:rq}. On the deferred part, the present finding
speaks in one direction only. As the Swiss case above already shows, hydrogen ranks low
among the competing uses for the molecule, because low-temperature space heating
measures it against a mature heat pump while hard-to-abate industry and parts of heavy
transport measure it against costly or absent alternatives
\citep{Quarton2019Curious, Brown2018Synergies}. The plausible spillover therefore runs from those uses back to buildings, where
scaling electrolysis and seasonal storage for industry and power may lower the
cost of the narrow building niche identified in
Section~\ref{sec:results:meritorder} without widening it. A full treatment, in
which buildings, industry, transport and power compete for a common hydrogen
supply under a single cost-minimising frame, is left to a subsequent paper.

%── 6.9 Limitations ────────────────────────────────────────────────────────────
\subsection{Limitations}

Three limitations bear on the discussion above, ahead of the fuller treatment in
Appendix~\ref{sec:limitations}. First, while the power layer now commits the firm fleet
with an explicit unit-commitment dispatch, the framework still does not embed a dynamic
capacity-expansion model. The mass deployment of heat pumps raises electricity demand materially; our own bridge puts study-area heat-pump electricity at 417~TWh in 2050 (Section~\ref{sec:power}), which calls for grid reinforcement and flexibility measures that the dispatch coupling here
quantifies but does not co-optimise. Second, the demand reconstruction operates at NUTS3
resolution, the finer spatial unit for which harmonised European data are available, with
three building-type categories, single-family, multi-family including apartments, and
other residential, trading intra-regional heterogeneity, in particular the dense
historic urban cores that might support a hydrogen niche, for cross-national
comparability across the 29 countries. Third, and now the most consequential of the
three, the non-fuel cost components of the hydrogen boiler, its capital and operating
costs, are less well established than those of mature technologies and carry wider
parameter uncertainty than the heuristic ranges applied here convey.

That uncertainty is close to a rounding term wherever fuel accounts for almost the whole
of the delivered cost. Here it is not, since more than two-fifths of the hydrogen
boiler's 2050 cost sits in capital, delivery and seasonal storage, the components
with the thinnest evidence base, and it is that cost which decides the country-level
base-load ranking. To those three limitations should be added the fiscal asymmetry set
out above, which runs in hydrogen's favour and so bounds the result only from the side
that would widen the heat pump's margin, and which Appendix~\ref{sec:limitations} takes
up in full. That appendix ranks the simplifications by consequence and leads with a
different four, the synthetic weather year behind every power-sector number, the two
network charges the symmetry does not cover, the same fiscal asymmetry, and the
energy-only market behind the recovery figures. The three named here bear on the
discussion above; the four named there bear on the results, and a reader weighing what
the paper can carry should read both lists.
